\hypertarget{results-and-discussion}{%
\chapter{Results and Discussion}\label{results-and-discussion}}

\begin{draftnote}

\textbf{Draft status, 14 November 2026.} Every figure below is traceable
to a committed script in \texttt{apps/detector/measurements/}. The one
figure that was not --- an AUC pair quoted in \S{}7.4 --- was retracted on
04 November and is marked as such in place rather than deleted.

\S{}7.8 reports the Layer 3 measurements, added 14 November. The layer is
measured but not deployed; \S{}5.4 states the difference.

This chapter reports a \textbf{negative result} for the project's
headline claim. That is deliberate and it is defensible, but \S{}7.10 must
be read with the project guide before the report is finalised.

\end{draftnote}

\begin{center}\rule{0.5\linewidth}{0.5pt}\end{center}

\hypertarget{what-was-measured-and-what-was-not}{%
\section{What was measured, and what was
not}\label{what-was-measured-and-what-was-not}}

Five measurement programmes exist, each a committed and re-runnable
script rather than a recorded number:

\begin{longtable}[]{@{}
  >{\raggedright\arraybackslash}p{(\columnwidth - 4\tabcolsep) * \real{0.3333}}
  >{\raggedright\arraybackslash}p{(\columnwidth - 4\tabcolsep) * \real{0.3333}}
  >{\raggedright\arraybackslash}p{(\columnwidth - 4\tabcolsep) * \real{0.3333}}@{}}
\toprule\noalign{}
\begin{minipage}[b]{\linewidth}\raggedright
programme
\end{minipage} & \begin{minipage}[b]{\linewidth}\raggedright
script
\end{minipage} & \begin{minipage}[b]{\linewidth}\raggedright
what it establishes
\end{minipage} \\
\midrule\noalign{}
\endhead
\bottomrule\noalign{}
\endlastfoot
Authorship discrimination & \texttt{python\_authorship.py},
\texttt{java\_authorship.py} & whether the style features distinguish
authors at all \\
Structural distribution & \texttt{structural\_floor.py} & how Layer 1
scores honest work by different authors \\
Labelled outsourcing corpus & \texttt{style\_shift.py} & a corpus of
partially outsourced files to test against \\
The ablation and its follow-ups & \texttt{ablation.py},
\texttt{layer2\_direction.py}, \texttt{suppression.py},
\texttt{independence.py} & whether Layer 2 adds signal beyond Layer 1,
and why not \\
Machine-generated code & \texttt{collect\_human.py},
\texttt{ai\_corpus.py}, \texttt{check\_ai\_corpus.py},
\texttt{baseline\_classifier.py} & whether generated code is separable
from student code, and at which edit grades \\
\end{longtable}

\textbf{What is not measured, and is a standing limitation.} There is no
precision or recall figure for Layer 1 on \emph{known-copied} pairs,
because no labelled corpus of genuine student plagiarism exists. The
structural measurement characterises the distribution over honest work
by different authors, which is a different claim and is reported as
such. Manufacturing a labelled corpus of real plagiarism would be worse
than admitting its absence.

\hypertarget{methodology-and-how-the-measurements-were-validated}{%
\section{Methodology, and how the measurements were
validated}\label{methodology-and-how-the-measurements-were-validated}}

This section is unusual in an undergraduate report and it is the
strongest thing in the chapter, so it comes before the results rather
than after.

Three days of this project were spent finding errors in its own
measurements. A platform-dependent sort; a single arbitrary sample
quoted to three significant figures; an alias merge applied after the
ownership share instead of before; a lift metric that was never
comparable across runs with different author counts. Several published
claims were retracted as a result.

Two permanent features of the measurement scripts exist \emph{because}
of those errors, and that is the right way to present them:

\begin{itemize}
\tightlist
\item
  \textbf{Sampling ranges, not point estimates.} Every figure is
  reported across 25 draws, because a single sample proved unreliable.
\item
  \textbf{Fixed author counts in every comparison table.} Lift is
  accuracy over chance and chance is \texttt{1/k}, so lift rises with
  the number of authors even when nothing about the features changed:
  the same Python data gives \textbf{2.43x at 5 authors and 4.65x at
  20}. Any two lift figures quoted side by side must hold the author
  count fixed.
\end{itemize}

Four errors were found in the ablation study itself during November, all
by reading the code rather than running it, and each had produced output
that looked like a result. They are reported here because the pattern is
the finding:

\begin{enumerate}
\def\labelenumi{\arabic{enumi}.}
\tightlist
\item
  The positive control compared the \textbf{wrong file} --- the corpus
  recorded the donor's author but not the donor file, so the control
  substituted an unrelated file by that author. It would have printed
  ``Layer 1 misses verbatim copying'', a false statement about the
  project's own core layer.
\item
  The chance baseline for precision@k was computed as spliced files over
  all scored files. Because each draw takes one of two versions per
  pair, chance is \textbf{0.5 by construction}; the error put it at
  0.429 and turned a non-result into an apparent thirteen-point lift.
\item
  Layer 1 was compared against \textbf{every} candidate in the pool,
  while the deployed detector shortlists and compares at most ten
  (\texttt{MAX\_DEEP\_CANDIDATES}). Comparing exhaustively credits Layer
  1 with matches the API cannot make, which biases the study
  \emph{against} Layer 2.
\item
  The sign test's verdict was printed without a \textbf{direction}. A
  two-sided test fires on any departure from an even split, so
  ``significant'' appeared beside a result that was significant evidence
  \emph{against} the hypothesis.
\end{enumerate}

A fifth was methodological: the first stratification bucketed on the
\emph{requested} splice fraction when the corpus documentation states
the \emph{actual} fraction is what an analysis must use. Correcting it
did not change the conclusion, which is the only reason the earlier
figures survive.

\hypertarget{layer-1-structural-similarity-on-honest-work}{%
\section{Layer 1 --- structural similarity on honest
work}\label{layer-1-structural-similarity-on-honest-work}}

Measured across 15 Python repositories: 1,989 candidate files, 466
single-author, 204 selected across 8 repositories and 27 authors, 100
pairs per category at seed 0, under production comparison settings.

\textbf{Only 8 of the 15 repositories produced any qualifying author.}
That is a sampling fact the chapter states rather than omits, and it
bounds how far the distribution generalises.

\textbf{A retraction belongs here.} An earlier claim that Layer 1
similarity has a floor near 0.5 was retracted. The argument assumed two
trees of the same shape differing only in labels; real unrelated code
differs in shape as well, so the floor does not hold.

\hypertarget{layer-2-does-the-style-signal-distinguish-authors}{%
\section{Layer 2 --- does the style signal distinguish
authors?}\label{layer-2-does-the-style-signal-distinguish-authors}}

This is the prior question. If the features cannot tell authors apart at
all, nothing downstream can work.

Corpus: 327 files with \textbf{zero parse failures} --- Java 123 files,
15 authors, 11 repositories; Python 204 files, 27 authors, 15
repositories.

\textbf{The honest headline, stated the way the measurement states it:}
across eleven within-project tests, the features reach roughly
\textbf{1.4x to 2.5x over chance where they work at all, and they fail
to beat chance in more than half the projects tested} --- five of eleven
hold, six do not.

\begin{itemize}
\tightlist
\item
  Strongest: \textbf{sympy}, 9 authors, 68 files, \textbf{2.51x}, never
  below 2.12x across 25 draws.
\item
  Weakest that still matters: \textbf{scikit-learn}, 7 authors, 54
  files, which does not clear chance at all.
\end{itemize}

Sample size is therefore not the whole story: projects differ in how
much personal style survives their own conventions. Any cross-project
figure is inflated by the repository confound, which accounts for much
of it.

\textbf{The feature selection is retracted.} The four scored features
were chosen from a 60-file sample on Day 19. Measured over 204 Python
files from 15 repositories:

\begin{itemize}
\tightlist
\item
  \texttt{blank\_line\_ratio}, ranked \textbf{first at 2.73} in the
  original sample and the reason the four were chosen, ranks
  \textbf{ninth of ten at 1.08} in Python.
\item
  \textbf{Nine of ten} features clear F = 1 in Python, not four.
\item
  Using every applicable feature \textbf{beats} the fixed four in both
  languages: Python \textbf{2.91x to 5.56x}, Java \textbf{2.80x to
  3.88x}.
\end{itemize}

A fixed four cannot be right for both languages in any case:
\texttt{blank\_line\_ratio} carries almost no author signal in Python
(1.08) while doing real work in Java (3.34); \texttt{comment\_density}
is the reverse in degree (13.97 against 2.49).

\textbf{RETRACTED, 04 November 2026.} An AUC comparison of
\textbf{0.766} using a student's own spread against \textbf{0.909} using
the cohort's was quoted in earlier project documentation as the evidence
for this choice. It is withdrawn. No script in
\texttt{apps/detector/measurements/} computes an AUC or an ROC curve at
all --- the metrics the measurement programme computes are lift, chance,
F ratio, accuracy and precision --- so the figures are reproducible from
nothing in the repository. Their likeliest source is the deleted Day 19
scratch file, which is the same provenance as the 2.27x and 3.15x
figures already retracted for the same reason.

\textbf{The design choice stands on its argument, which is not a
measurement and should not be presented as one.} A student's own spread
is estimated from three or four anchors and is too unstable to divide
by; a mean is personal, whereas a spread is closer to a property of the
feature. Establishing the cohort denominator empirically would need a
run that does not yet exist.

\hypertarget{the-ablation-does-layer-2-add-signal-beyond-layer-1}{%
\section{The ablation --- does Layer 2 add signal beyond Layer
1?}\label{the-ablation-does-layer-2-add-signal-beyond-layer-1}}

\hypertarget{the-corpus}{%
\subsection{The corpus}\label{the-corpus}}

\texttt{style\_shift.py} builds the labelled set the earlier
measurements lacked. It takes a file one author owns and replaces some
of its top-level functions with another author's, which is what
outsourcing part of an assignment looks like from outside: the
submission is mostly yours and a section of it is not. Both versions of
every file are written, sharing a \texttt{pair\_id}.

From 466 single-author files, 204 kept after the per-author cap, 27
authors: \textbf{214 pairs written, 155 splices discarded (42\%), 24
authors represented.}

\begin{longtable}[]{@{}lll@{}}
\toprule\noalign{}
requested fraction & pairs & median actual \\
\midrule\noalign{}
\endhead
\bottomrule\noalign{}
\endlastfoot
0.10 & 61 & 0.104 \\
0.25 & 80 & 0.235 \\
0.50 & 73 & 0.438 \\
\end{longtable}

Three corpus properties the chapter must state:

\begin{itemize}
\tightlist
\item
  \textbf{The 42\% discard rate is the honest figure, not a failure.} A
  splice is discarded when the file is too structurally coarse to hit
  the requested fraction, when no donor unit matches the size of what it
  would replace, or when the result no longer parses. Keeping them would
  mean a bucket labelled 0.10 containing splices at 0.40.
\item
  \textbf{Three of 27 authors produced no usable pair}, so the corpus is
  biased toward authors whose files are structurally amenable to
  splicing.
\item
  \textbf{214 pairs come from only 94 distinct source files.} Eighty-one
  sources were reused across the three fractions, producing
  byte-identical clean copies. The statistics are therefore
  \textbf{stratified by fraction}, within which each source appears at
  most once. \textbf{Ninety-four, not 214, is the number every figure
  below rests on.}
\end{itemize}

The length confound was controlled by size-matching donor units to
within 30\% of what they replace. The control worked: median length
delta is \textbf{-1 line}, and the measured correlation between absolute
line delta and the RPS delta is \textbf{-0.131}, under 2\% of variance
and negative. Detection is not tracking file size.

\hypertarget{the-positive-control-which-passes}{%
\subsection{The positive control, which
passes}\label{the-positive-control-which-passes}}

With the donor's own file inside Layer 1's pool --- a verbatim copy the
system can see --- Layer 1 behaves exactly as it should:

\begin{longtable}[]{@{}llll@{}}
\toprule\noalign{}
splice fraction & Layer 1 gap & paired sign test & precision@10 \\
\midrule\noalign{}
\endhead
\bottomrule\noalign{}
\endlastfoot
0.10 & +0.0150 & 42 of 60, p = 2.67 \texttimes{} 10\textsuperscript{-3} & 0.700 \\
0.25 & +0.0335 & 53 of 80, p = 4.87 \texttimes{} 10\textsuperscript{-3} & 0.800 \\
0.50 & +0.0706 & 47 of 73, p = 1.86 \texttimes{} 10\textsuperscript{-2} & \textbf{1.000} \\
\end{longtable}

Monotonic in how much was copied, significant at every fraction, and
perfect ranking at the top ten when half the file is someone else's.
Layer 2's figures are byte-identical between the two conditions, as they
must be, since the donor never enters Layer 2's computation.

\textbf{This is what licenses the negative result below.} The apparatus
demonstrably detects what it is supposed to detect.

\hypertarget{the-result-donor-excluded-the-actual-thesis-test}{%
\subsection{The result, donor excluded --- the actual thesis
test}\label{the-result-donor-excluded-the-actual-thesis-test}}

This is the threat model: outsourced work with no second copy in the
cohort.

\begin{longtable}[]{@{}
  >{\raggedright\arraybackslash}p{(\columnwidth - 8\tabcolsep) * \real{0.2000}}
  >{\raggedright\arraybackslash}p{(\columnwidth - 8\tabcolsep) * \real{0.2000}}
  >{\raggedright\arraybackslash}p{(\columnwidth - 8\tabcolsep) * \real{0.2000}}
  >{\raggedright\arraybackslash}p{(\columnwidth - 8\tabcolsep) * \real{0.2000}}
  >{\raggedright\arraybackslash}p{(\columnwidth - 8\tabcolsep) * \real{0.2000}}@{}}
\toprule\noalign{}
\begin{minipage}[b]{\linewidth}\raggedright
splice fraction
\end{minipage} & \begin{minipage}[b]{\linewidth}\raggedright
Layer 1 gap
\end{minipage} & \begin{minipage}[b]{\linewidth}\raggedright
Layer 2 gap
\end{minipage} & \begin{minipage}[b]{\linewidth}\raggedright
precision@10, L1
\end{minipage} & \begin{minipage}[b]{\linewidth}\raggedright
precision@10, L1+L2
\end{minipage} \\
\midrule\noalign{}
\endhead
\bottomrule\noalign{}
\endlastfoot
0.10 & +0.0095 & \textbf{-0.0124} & 0.500 & 0.500 \\
0.25 & +0.0088 & \textbf{-0.0176} & 0.700 & \textbf{0.400} \\
0.50 & -0.0006 & \textbf{-0.0090} & 0.700 & 0.500 \\
\end{longtable}

Chance is 0.500 by construction.

Three statements, all negative:

\begin{enumerate}
\def\labelenumi{\arabic{enumi}.}
\tightlist
\item
  Layer 2 scores \textbf{lower} on spliced files than on clean ones at
  every fraction. That is the wrong direction, not merely a small
  effect.
\item
  Adding Layer 2 to Layer 1 \textbf{degrades} ranking precision at every
  fraction, falling below chance at 0.25.
\item
  There is \textbf{no dose-response}. Measured continuously against the
  actual splice fraction rather than a bucket: \textbf{+0.022} for Layer
  2.
\end{enumerate}

\hypertarget{why-layer-2-fails-and-it-is-structural}{%
\section{Why Layer 2 fails, and it is
structural}\label{why-layer-2-fails-and-it-is-structural}}

Two diagnostics separate the possible causes.

\hypertarget{the-features-do-move}{%
\subsection{The features do move}\label{the-features-do-move}}

\texttt{suppression.py} measures how far there was to move against how
much of that move happened. \textbf{Reach} is the donor-to-target
distance in cohort standard deviations; \textbf{realisation} is the
fraction of the predicted displacement that actually occurred.

\begin{longtable}[]{@{}lll@{}}
\toprule\noalign{}
feature & median \textbar reach\textbar{} & median realisation \\
\midrule\noalign{}
\endhead
\bottomrule\noalign{}
\endlastfoot
\texttt{blank\_line\_ratio} & 0.945 & \textbf{0.997} \\
\texttt{avg\_line\_length} & 0.905 & \textbf{0.907} \\
\texttt{comment\_density} & 0.590 & \textbf{0.893} \\
\texttt{max\_block\_depth} & 0.694 & \textbf{-0.000} \\
\end{longtable}

Predicted against observed displacement, in standard deviations:

\begin{longtable}[]{@{}lll@{}}
\toprule\noalign{}
requested fraction & median \textbar predicted\textbar{} & median
\textbar observed\textbar{} \\
\midrule\noalign{}
\endhead
\bottomrule\noalign{}
\endlastfoot
0.10 & 0.085 & 0.062 \\
0.25 & 0.182 & 0.254 \\
0.50 & 0.365 & 0.443 \\
\end{longtable}

So the donors genuinely differ from the targets, and for three of the
four features the simple mixing model holds almost exactly. \textbf{The
corpus is not suppressing the signal.} The fault is in the feature set
and the score.

\hypertarget{an-rms-distance-is-direction-blind}{%
\subsection{An rms distance is
direction-blind}\label{an-rms-distance-is-direction-blind}}

The score is \texttt{min(1,\ rms(z)\ /\ Z\_SATURATION)} --- a
\textbf{magnitude} of deviation from the student's own baseline. The
splice displaces a file by about 0.44 standard deviations per feature at
the half-replaced fraction. But the direction of that displacement is
set by where the donor happens to sit relative to the student, which is
arbitrary across pairs.

A point already 0.55 to 0.90 standard deviations from its baseline,
displaced 0.44 in an arbitrary direction, moves \textbf{closer} to the
baseline about as often as further from it. The measurement confirms it:
at the half-replaced fraction, \textbar z\textbar{} grew in \textbf{37
of 73} pairs for \texttt{blank\_line\_ratio}, \textbf{38 of 73} for
\texttt{avg\_line\_length} and \textbf{37 of 68} for
\texttt{comment\_density}. Coin flips. That is why the gaps in \S{}7.5.3
are slightly negative rather than merely small.

\texttt{Z\_SATURATION\ =\ 3} clamps \textbf{zero} pairs, so saturation
hides nothing; the layer's usable range on this corpus is roughly 0.18
to 0.30 out of nought to one.

\hypertarget{a-maximum-is-not-a-mixable-statistic}{%
\subsection{A maximum is not a mixable
statistic}\label{a-maximum-is-not-a-mixable-statistic}}

\texttt{max\_block\_depth} has reach 0.694 --- donors do differ on it
--- yet realisation is exactly \textbf{-0.000}: the splice usually does
not change it at all. Replacing some functions cannot lower a maximum
unless the deepest is replaced, and a donor's function raises it only if
it is deeper. It is tied in \textbf{57 of 61} pairs at the 0.10
fraction. Its realisation upper quartile is 2.250, so when it does move
it overshoots wildly.

Neither failure is a tuning problem. Both say the same thing: \textbf{a
whole-file magnitude cannot detect a partial, signed substitution.}

\hypertarget{the-arbitrary-direction-measured-directly}{%
\subsection{The arbitrary direction, measured
directly}\label{the-arbitrary-direction-measured-directly}}

The claim above --- that the donor's position relative to the student is
arbitrary across pairs --- was initially inferred from the coin-flip
counts in \S{}7.6.2. \texttt{suppression.py} now measures it, reporting the
\textbf{sign} of reach and not only its magnitude:

\begin{longtable}[]{@{}lllll@{}}
\toprule\noalign{}
feature & donor above target & of & share & p \\
\midrule\noalign{}
\endhead
\bottomrule\noalign{}
\endlastfoot
\texttt{blank\_line\_ratio} & 102 & 214 & 0.477 & 0.54 \\
\texttt{avg\_line\_length} & 107 & 214 & \textbf{0.500} & 1.00 \\
\texttt{max\_block\_depth} & 85 & 174 & 0.489 & 0.82 \\
\texttt{comment\_density} & 96 & 208 & 0.462 & 0.30 \\
\end{longtable}

All four are indistinguishable from a coin flip, and
\texttt{avg\_line\_length} is exactly half. The displacement a splice
applies therefore has a magnitude the mixing model predicts and
\textbf{a sign that carries no information} --- which is precisely the
condition under which a distance from a baseline cannot detect it.

The account in \S{}7.6 is consequently measured at every link rather than
argued at any of them: the donors differ (reach), the direction of that
difference is arbitrary (this table), and the splice realises the
predicted displacement (realisation). A magnitude applied in an
arbitrary direction to a point already off centre moves it closer about
as often as further away, and the paired counts in \S{}7.6.2 show exactly
that.

\hypertarget{layer-independence}{%
\section{Layer independence}\label{layer-independence}}

The RPS adds the layers with fixed weights, which presumes each says
something the other does not. Measured from the saved per-item scores:
Pearson \textbf{-0.042} across all 428 scored files, about -0.04 within
each class, Spearman within \textpm{}0.11 everywhere.

The weighting assumption therefore holds. It should be read carefully
rather than as a vindication: \textbf{two signals are also uncorrelated
when one of them is noise.} Independence is necessary for the
composition to make sense; it is not by itself evidence that combining
the layers helps. On this corpus, \S{}7.5.3 shows it does not.

\hypertarget{layer-3-machine-generated-code-measured}{%
\section{Layer 3 --- machine-generated code,
measured}\label{layer-3-machine-generated-code-measured}}

Chapter 5 \S{}5.4 records why Layer 3 went unbuilt for most of this
project: the published approach depends on training data that resembles
the submissions it will police, and no such corpus existed. That
objection was answered by building the corpus rather than by assuming
one, and the layer has now been measured. It is \textbf{not deployed}
--- \S{}5.4 explains what is and is not wired into the detection path ---
but the question of whether it works is no longer open.

\hypertarget{the-corpus-and-one-serious-false-start}{%
\subsection{The corpus, and one serious false
start}\label{the-corpus-and-one-serious-false-start}}

The first attempt reused the 214-pair splice corpus from \S{}7.5. That
corpus is CPython standard library code, which is appropriate for asking
whether a splice is detectable and wholly inappropriate here. A
classifier given library code as its human class and generated code as
its AI class separates twenty years of committee-edited library source
from freshly written single-purpose modules. It would have scored well
and measured nothing about authorship. The median file length alone
makes the mismatch plain: \textbf{444 lines for the standard library
against 98 for student coursework}.

The human class was rebuilt from scratch. \texttt{collect\_human.py}
clones public GitHub repositories tagged as coursework, mini-projects
and student projects, and keeps a file only if \textbf{that file's own
last commit} predates 1 January 2023. Repository age is not sufficient
--- a 2019 repository can hold a file edited in 2024 --- so the filter
asks git for each candidate individually. The result is 96 files from 25
repositories, last-commit dates spanning 2017-02-28 to 2022-11-30, with
per-file provenance recorded in a manifest.

The AI class is 60 files generated to match the human corpus on two
axes.

\textbf{Task mix}, so the classifier cannot separate the classes by
subject matter:

\begin{longtable}[]{@{}ll@{}}
\toprule\noalign{}
share & task type \\
\midrule\noalign{}
\endhead
\bottomrule\noalign{}
\endlastfoot
28\% & computer vision and machine-learning coursework \\
23\% & networking, bots and scraping \\
21\% & tools, editors, calculators \\
14\% & CRUD systems with a GUI and a database \\
7\% & games \\
7\% & algorithms \\
\end{longtable}

\textbf{Surface style}, which matters more than it appears to. Language
models produce docstrings and type hints on almost everything, while the
human corpus carries module docstrings on 11\% of files, function
docstrings on 33\% and type hints on 11\%. Left alone, those four traits
separate the classes almost perfectly and a classifier learns formatting
rather than authorship. Each prompt therefore carried an explicit style
instruction drawn at the human corpus's own rates, and
\texttt{check\_ai\_corpus.py} verified every generated file against what
its prompt asked for. Two files in fifty ignored the instruction and
were corrected; models are not reliable about this.

\hypertarget{matched-preprocessing-and-why-the-control-is-not-a-deduction}{%
\subsection{Matched preprocessing, and why the control is not a
deduction}\label{matched-preprocessing-and-why-the-control-is-not-a-deduction}}

\texttt{ai\_corpus.py} builds three edit grades: \textbf{raw} model
output, \textbf{disguised} output with identifiers renamed and comments
and docstrings stripped, and \textbf{mixed}, one generated function
spliced into an otherwise human file at a minimum AI share of 0.10.

Grade 2 carries a trap. If only the AI class is disguised, the stripped
comments and \texttt{v0,\ v1,\ v2} identifiers appear in exactly one
class by construction, and the classifier wins by detecting the
transformation. On an early run with deliberately falsified labels this
produced \textbf{100\% recall on grade 2 that meant nothing at all}.

The fix is matched preprocessing: a disguised copy of every human file
is added to the human class, so the transformation appears on both sides
and cannot act as a class cue. With matching in place and the labels
still falsified, grade-2 recall fell to \textbf{0\%} --- the correct
answer when there is nothing to find, and the negative control the rest
of this section rests on.

A control measurement accompanies every run, training on plain human
code against disguised human code. Both classes are human-written, so
the figure is purely how detectable the transformation is in isolation:
\textbf{0.8543}. An earlier version of the measurement script instructed
the reader to subtract that from the grade-2 result. \textbf{That
instruction was wrong and has been removed.} Once preprocessing is
matched the disguise sits in both classes and cannot inflate anything;
the control states why matching was necessary, not how much of the
result is artifact.

\hypertarget{the-result}{%
\subsection{The result}\label{the-result}}

351 items in 102 groups, split 233 train and 118 test. Items derived
from one source file are the same program in three disguises, and a
mixed file is mostly its human host, so whole groups are held out rather
than individual files.

\begin{longtable}[]{@{}ll@{}}
\toprule\noalign{}
measure & value \\
\midrule\noalign{}
\endhead
\bottomrule\noalign{}
\endlastfoot
ROC AUC, overall & \textbf{0.9521} \\
specificity --- human files left unflagged & 41 / 44 (93.2\%) \\
recall, grade 1 --- raw generated output & 22 / 25 (88.0\%) \\
recall, grade 2 --- disguised & 20 / 25 (80.0\%) \\
recall, grade 3 --- one function spliced in & \textbf{9 / 24
(37.5\%)} \\
\end{longtable}

\textbf{The grade-3 figure is the result.} Wholesale generation is
caught most of the time; partial insertion is caught barely more often
than a coin would manage. The classifier sees a file as one bag of
n-grams, so a single generated function is diluted by the human host
surrounding it.

That is the same failure, reached independently, as \S{}7.5.3 and \S{}7.6:
Layer 2 could not see a partial substitution either, for reasons that
were structural rather than fixable. \textbf{Two layers, two methods,
two corpora, one conclusion} --- this class of detector catches
wholesale substitution and misses partial substitution. The agreement is
worth more than the headline figure, because it is a statement about the
problem rather than a score for an implementation.

\hypertarget{two-confounds-tested-and-rejected}{%
\subsection{Two confounds tested and
rejected}\label{two-confounds-tested-and-rejected}}

A figure of 0.95 invites the question of what else it could be
measuring. Two candidates were checked.

\textbf{Era.} The human corpus is 2017--2022 and the AI corpus is 2026.
If Python idiom moved in that window, the classifier could be separating
dates rather than authors. Measured across 47 human and 60 generated
files:

\begin{longtable}[]{@{}lll@{}}
\toprule\noalign{}
marker & human & generated \\
\midrule\noalign{}
\endhead
\bottomrule\noalign{}
\endlastfoot
f-strings & 17\% & 25\% \\
type hints & 9\% & 12\% \\
\texttt{pathlib} & 2\% & 2\% \\
walrus \texttt{:=} & 0\% & 0\% \\
\end{longtable}

No meaningful gap. The students wrote structurally modern Python, and
this explanation is rejected.

\textbf{Boilerplate.} \texttt{if\ \_\_name\_\_\ ==\ "\_\_main\_\_"}
appears in 82\% of AI items and 32\% of human items, and reaches
\textbf{AUC 0.7499 on its own} --- a single line, trivially added or
deleted by anyone wishing to evade detection. The full model scores
0.9521 with it present and \textbf{0.9616 with it removed entirely}.
With 118 test items the standard error on an AUC is roughly 0.025, so
those two figures are indistinguishable: the feature contributes nothing
either way, and the result does not rest on it.

\hypertarget{what-is-not-claimed}{%
\subsection{What is not claimed}\label{what-is-not-claimed}}

\begin{itemize}
\tightlist
\item
  \textbf{CodeBERT is untrained.} \texttt{train\_codebert.py} exists,
  its non-model components are unit tested, and its training loop has
  never executed. Every figure above belongs to the logistic-regression
  baseline. \textbf{No CodeBERT accuracy is claimed anywhere in this
  report.}
\item
  \textbf{The generated class is dominated by one model.} 40 of 60 files
  came from one assistant and 20 from others. Generalisation to other
  generators is untested.
\item
  \textbf{The generated files use \texttt{\%} string formatting far more
  than is natural} --- 19\% of human files against 55\% of generated
  ones --- because the author chose that style while producing them, to
  avoid an obviously generated look. This is not natural model output.
  Its direction makes the task harder rather than easier, so it does not
  inflate the result, but it is an artifact introduced by the corpus
  author and is disclosed here rather than left to be found.
\item
  \textbf{Grade-3 hosts come from a narrow pool.} Only 21 of 96 human
  files have enough top-level functions to host a splice, so 54 mixed
  items draw on 21 hosts.
\item
  \textbf{The corpus is small.} 351 items, 102 groups, and each
  per-grade recall rests on 24 or 25 test items. Every figure carries
  wide error bars.
\item
  \textbf{Python only}, as with the ablation in \S{}7.5.
\item
  \textbf{One file is unaccounted for.} The collector wrote 97 files and
  the corpus builder found 96 usable, under what should be identical
  filters. The discrepancy is unexplained and is recorded rather than
  quietly absorbed.
\end{itemize}

\hypertarget{limitations}{%
\section{Limitations}\label{limitations}}

\begin{enumerate}
\def\labelenumi{\arabic{enumi}.}
\tightlist
\item
  \textbf{No labelled corpus of genuine plagiarism.} All structural
  figures describe honest work by different authors.
\item
  \textbf{94 independent units, not 214.} Stated at every figure in
  \S{}7.5.
\item
  \textbf{Open-source library code, not coursework.} Repository
  conventions are a confound, and the within-project figures are the
  defensible ones.
\item
  \textbf{Corpus biased toward splice-amenable files}; 3 of 27 authors
  produced none.
\item
  \textbf{One language.} The ablation ran on Python only. The authorship
  measurement covers both, and the two features most affected behave
  oppositely across them, so the ablation should not be assumed to
  transfer to Java.
\item
  \textbf{Layer 3 measured but not deployed.} The renormalisation is
  always active in the running system, so the three-layer design remains
  untested as a three-layer design even though Layer 3 now has figures
  of its own (\S{}7.8).
\item
  \textbf{Discrimination is not detection.} \S{}7.4 shows the features
  carry some author signal; \S{}7.5 shows that signal does not survive
  partial substitution. The two are not in conflict --- identifying a
  whole file's author is a different task from noticing that a fraction
  of it changed hands --- and the report must keep them separate.
\end{enumerate}

\hypertarget{what-this-means-for-the-projects-claim}{%
\section{What this means for the project's
claim}\label{what-this-means-for-the-projects-claim}}

Layer 2 is the only layer that can see outsourced work with no second
copy anywhere. That capability is what distinguishes this project from a
clone detector, and \textbf{the one corpus built to test it does not
support it.}

The honest statement of the contribution is therefore narrower than the
one the project set out with, and it is still a real contribution:

\begin{itemize}
\tightlist
\item
  a working two-layer evidence system with renormalised, auditable
  scoring and stored evidence;
\item
  five baseline-integrity mitigations that make the system report
  \emph{that it cannot tell} rather than emit a number that looks like
  an answer;
\item
  a reproducible measurement programme that establishes where the
  behavioural approach works (author discrimination, within project,
  1.4x to 2.5x where it works at all) and where it does not (partial
  outsourcing, not at all);
\item
  a diagnosis of \emph{why} it does not, precise enough to name the fix:
  per-unit rather than whole-file features, and a signed rather than
  magnitude comparison, so that moving \emph{toward} another author's
  typical values is itself the signal.
\end{itemize}

That last item is a design change and not a reweighting, and it is
future work.

\begin{draftnote}

\textbf{Action required before the report is finalised.} The project's
claimed research novelty, as stated in the presentation and in Chapter
1, asserts behavioural detection of outsourced work. \S{}7.5 does not
support that assertion. \textbf{This must be reworded in consultation
with the project guide, Mr.~Ashish Sharma, and not unilaterally.}
Chapter 2 \S{}2.7 carries the same flag. A negative result with a mechanism
and a named fix is a defensible final-year outcome; an overclaim
contradicted by the project's own committed measurement script is not.

\end{draftnote}
